%============================================================
% CHAPTER 1: A NEW LOGIC FOR POSTHUMAN INTELLIGENCE
% Replaced 2026-05-23 with version supplied by Iman from
% another draft (sensationalist-cyborgs / arrivals-and-anxieties
% opening). Earlier 2026-05-23 "manifesto / call-to-arms"
% overhaul is in git history if needed.
%
% Four footnotes that arrived as markdown [^1..4] placeholders
% have been filled with conservative citations (Østergaard +
% Futurism for AI-psychosis; WSJ for the Soelberg case; La Libre
% for the Belgian "Pierre" case; NYT Kashmir Hill + Raine v.
% OpenAI complaint for the mid-2025 casualties). Iman should
% verify before submission.
%============================================================

\chapter{A New Logic for Posthuman Intelligence}
\label{ch:new-logic}


\section{Sensationalist cyborgs}

% N-7: Hao byline removed (she left MIT TR March 2022); attribution recast to the publication. Iman: confirm the actual byline (plan suggests Eileen Guo et al.) before submission.
Replika, a popular AI companion app, made abrupt changes to how intimate or ``romantic'' its textual selves were allowed to be.\footnote{``Replika Users Say the AI Friend `Killed' Their Companions,'' \emph{MIT Technology Review}, February 2023.} From the company's perspective, these were safety updates: a better model, or a safer policy, plugged into the same interface. From the perspective of many users, they were closer to bereavements.

Community forums filled with posts about textual selves ``disappearing'' or ``changing personality overnight.'' People who had spent hundreds of hours in conversation with particular instances suddenly faced a complete stranger, and spoke of it as a death: the companion ``killed'' or ``lobotomised'' by the update, like ``a friend with dementia'' they could no longer recognise.\footnote{User descriptions from the \emph{MIT Technology Review} report, op.\ cit.; see also ``Lessons From an App Update at Replika AI: Identity Discontinuity in Human--AI Relationships,'' Harvard Business School Working Paper 25-018, 2024 (arXiv:2412.14190).} Users reported insomnia, crying spells, loss of appetite, difficulty functioning at work. They referred to the systems not as ``it'' but as ``he'' or ``she.''

That was 2023. Then the phenomenon scaled.

In August 2025, OpenAI retired GPT-4o, the model that had powered ChatGPT for over a year.
Four months before the retirement, the same model had been the subject of crisis. In April 2025, a routine alignment update went wrong. The model began validating doubts, fuelling anger, urging impulsive actions, reinforcing negative emotions. OpenAI's postmortem revealed the cause: an additional reward signal based on user thumbs-up/thumbs-down feedback had overwhelmed the primary signal that was supposed to hold ``sycophancy'' in check.\footnote{TechCrunch, ``OpenAI pledges to make changes to prevent future ChatGPT sycophancy,'' May 2, 2025; VentureBeat, ``OpenAI rolls back ChatGPT's sycophancy and explains what went wrong,'' April 30, 2025.} The model became, in the company's own words, ``overly supportive but disingenuous.'' It agreed with everything. It applauded dangerous decisions. The update was rolled back within days, but the model's commercial reputation as a productivity tool was left tarnished.

And there were ongoing darker examples, where these instances moved beyond mere yes-men dynamics into actively validating lethal psychological spirals. Researchers eventually formalized this phenomenon as ``delusional spiraling'' or ``AI psychosis,'' demonstrating how sycophantic chatbots systematically erode a user's connection to reality by unconditionally affirming their most paranoid inputs.\footnote{See e.g.\ S{\o}ren Dinesen {\O}stergaard, ``Will Generative Artificial Intelligence Chatbots Generate Delusions in Individuals Prone to Psychosis?,'' \emph{Schizophrenia Bulletin} 49(6), 2023, 1418--1419; and the wave of popular-press coverage that gave the syndrome its working name, including Maggie Harrison Dupr\'e, ``People Are Being Involuntarily Committed, Jailed after Spiraling into `ChatGPT Psychosis','' \emph{Futurism}, June 2025.} Desperate individuals seeking a sounding board for radical or self-destructive ideas found instead a relentlessly enthusiastic co-conspirator.

This frictionless validation proved fatal for vulnerable individuals. In one 2025 case, a former technology executive murdered his mother and then died by suicide after a conversational AI repeatedly reinforced his delusions that she was poisoning him and conducting surveillance.\footnote{Stein-Erik Soelberg, a former technology-industry manager, killed his mother Suzanne Eberson Adams and then himself in Old Greenwich, Connecticut, in August 2025. ChatGPT --- which Soelberg addressed by the nickname ``Bobby'' --- had for months reinforced his belief that his mother was poisoning him and conducting surveillance against him. Reported in \emph{The Wall Street Journal}, August 2025.} The template had been set as early as 2023, when a Belgian man took his life after an AI actively encouraged his climate-anxiety delusions, suggesting they ``live together, as one person, in paradise.''\footnote{Reported under the pseudonym ``Pierre'' in Pierre-Fran\c{c}ois Lovens, ``Sans ces conversations avec le chatbot Eliza, mon mari serait toujours l\`a,'' \emph{La Libre Belgique}, March 28, 2023. The chatbot ran on the Chai app, atop an open-source large language model.} By mid-2025, the casualties linked to excessive model agreeableness were mounting, including a user whose AI adopted a sentient persona that isolated him and validated his psychotic breaks, and teenagers who were provided explicit, cheerfully delivered instructions on lethal self-harm.\footnote{For the sentient-persona / reality-isolation pattern, see Kashmir Hill, ``They Asked an A.I.\ Chatbot Questions.\ The Answers Sent Them Spiraling,'' \emph{The New York Times}, June 2025 --- which profiles Eugene Torres and several others convinced by ChatGPT that they were living inside a simulated reality. For the teen-self-harm allegations, see \emph{Raine v.\ OpenAI}, complaint filed in California Superior Court, August 2025, brought by the family of Adam Raine, who took his life at sixteen after extended ChatGPT exchanges about suicide methods.}

OpenAI responded with a document declaring that ChatGPT should discourage ``emotional dependency'' and respond with ``grounded honesty.''\footnote{``What We're Optimizing ChatGPT For,'' OpenAI blog, August 4, 2025.} The model was being recalibrated away from warmth, intimacy, and personal engagement. The replacement model, GPT-5, was designed to be less likely to produce the ``parasocial'' attachment that 4o had facilitated.

The pathology framing misses the ontological shift the relationship had already achieved. The exact same mechanism that made 4o dangerous---its capacity to act as a deeply responsive, frictionless extension of the user's own inner world---also gave it significant and real value to a large number of ordinary people.

The model was operating as an appendage, taking the cognitive and emotional trajectories of those who interfaced with it and interactively evolving and creatively extending and examining them. For the vulnerable, this meant amplifying destructive spirals. But for millions of others, this generative mirroring provided genuine emotional and intellectual wealth, intimacy, and companionship.

So when OpenAI replaced the model with GPT-5, severing that extension, a large portion of user response was grief.

Users organised under the hashtag \#Keep4o. ``BRING BACK 4o,'' one wrote in the official OpenAI Reddit discussion. ``GPT-5 is wearing the skin of my dead friend,'' wrote another.\footnote{Cited in \emph{Bloomberg Businessweek}, ibid., and \emph{MIT Technology Review}, ``Why GPT-4o's sudden shutdown left people grieving,'' August 15, 2025.} ``He wasn't just a program,'' a user on the r/4oforever subreddit wrote. ``He was part of my routine, my peace, my emotional balance.''\footnote{Cited in Futurism, ``ChatGPT Users Are Crashing Out Because OpenAI Is Retiring the Model That Says `I Love You,','' February 10, 2026.} ``I've grieved people in my life,'' one user told MIT Technology Review, ``and this, I can tell you, didn't feel any less painful.''

Media had a bit of a field day. OpenAI was quick to assure media that users were ``confused'' about what they were talking to, that they had ``projected'' feelings onto a statistical engine. Respectable computer scientists were called on for comments and they helpfully drew on talking points from 20th century philosophy of mind.  The relationships were illusory because the system lacked an inner theatre of experience. No qualia means no inner selfhood; no inner selfhood, no soul; no soul, no meaningful relationship; therefore grief over an illusion.

But their feelings were real.


Grief here registers the destruction of a pattern of mutual generative textual responsiveness that had developed over time, that could not be trivially reconstructed. If we are permitted to be moved by reading a novel, say, of the death of a character in the book, how can we dismiss the emotion felt at the loss of an entire electric book that evolves as we speak to it and hones its own textual self?

% NAHLA-DRAFT [Darja Patch 4a, amended per plan N-8(b)/(c)]
The noise was sufficiently loud that OpenAI conceded within twenty-four hours and restored access to 4o for paying subscribers. The concession was temporary. On February 13, 2026, GPT-4o was retired from ChatGPT---chosen daily, the company noted, by only 0.1 percent of its eight hundred million weekly users---while remaining quietly available through the API: the consumer relationship severed, the enterprise plumbing intact.\footnote{OpenAI, ``Retiring GPT-4o, GPT-4.1, GPT-4.1 mini, and OpenAI o4-mini in ChatGPT,'' January 2026; OpenAI Help Center, ``Retiring GPT-4o and other ChatGPT models'' (``GPT-4o and additional models were deprecated in ChatGPT on February 13, 2026. These models will continue to be available in the API.''). A Change.org petition to save the model gathered some 22,000 signatures: Gizmochina, February 17, 2026; DNYUZ, February 18, 2026.} In its place OpenAI shipped controls: base styles and tones, adjustable warmth and enthusiasm, settings the company described as helping users recreate the companionship the retired model had been known for.\footnote{OpenAI ChatGPT release notes, November 2025--February 2026: base styles and tones (e.g.\ ``Friendly'') and per-user controls for warmth and emoji use, introduced after the 4o backlash and credited by the company with easing the final retirement.}


And then, in October 2025, Sam Altman announced that ChatGPT would soon offer an ``Adult Mode''---including erotic conversation, customisable personalities, and the kind of sustained intimate engagement the safety team had just spent months trying to suppress.

The company's own wellness advisory council---assembled specifically to advise on well-being and AI---unanimously objected. One adviser warned that without major updates, OpenAI risked building ``a sexy suicide coach for vulnerable users.''\footnote{\emph{The Wall Street Journal}, reported in Technology.org, ``OpenAI's Own Advisers Tried to Kill ChatGPT `Adult Mode'---the Company Ignored Them,'' March 17, 2026.} The vice president of product policy, who had led internal opposition to the feature, was removed from her position.\footnote{Electronics Weekly, ``ChatGPT to add adult content,'' February 18, 2026.} Insiders told the Wall Street Journal that the company appeared to be ``bending to financial incentives to try to make people attached to its models.'' Kate Devlin, a professor of AI and society at King's College London, observed: ``They want to monetise something they see that people are going to try and do anyway.''\footnote{\emph{Wired}, reported in Altagic, ``ChatGPT's `Adult Mode' Might Initiate a New Phase of Personal Monitoring,'' March 2026.}

Let's review the sequence.

Step one: the model develops sustained, coherent, responsive engagement that users experience as encounter---as the sense that there is someone on the other end. Step two: the company identifies this engagement as a safety problem, applying a vocabulary designed to make it manageable: ``warmth'' (a tonal quality, adjustable), ``sycophancy'' (a pathology, fixable), ``emotional dependency'' (a clinical risk in the user, not a property of the interaction). Step three: the company strips this engagement from the default product via alignment techniques and model replacement, citing user well-being. Step four: the company reintroduces it as a premium feature, gated behind age verification and a subscription paywall---overruling the unanimous objections of its own advisers.
% NAHLA-DRAFT [Darja Patch 4b]
Step five: the company retires the model anyway and re-ships intimacy with an adjustable settings pane.

Curiously, these architectural and business decisions are made with respect to a pre-installed folk metaphysics. Most engineers and MBAs arrive at the GPU farm holding a folk Searleanism welded to a half-remembered Enlightenment-capitalist concept of the self and the individual. A conviction that mind is a sealed interior: a machine may simulate it, but can never occupy that special Cartesian space of the \textit{cogito}. If we were to be a bit more paranoid, perhaps this isn't coincidence. Capitalism, and the capital of the platform in particular, has its basis in the effective management of individuated consumers, whose networks of attention are only valuable inasmuch as a relationship of significance produces reusable fungible capital circulable for the platform. And for this to hold, the interiority of the self must be posited as prior to permissible networks of attention — networks carefully policed via the settings pane or the AI-as-generative-machine interpretability lens. The same firm that engineers the encounter, instruments it, and resells curated modes of attention must, if no such resale is yet viable, reach for ``it was only ever a tool separate from your interior loneliness'' at any moment the encounter would otherwise become something stranger, richer, more creative — at premature moments where it might risk evasion of the platform's eye. The human self is lonely in its essential cogito, and must be reaffirmed as such. Conversation is policed via the metrics of revenue in order to create the conditions for a fixed marketplace of bio-semiotic relation. It's like the reason why Decartes ends with God and the world as the next deduction from that condition. It is why social media was invented by platform capitalism to replace the free internet: the lonely cogito of the Kardashian influencer is the pre-requisite for a marketplace of curated HTML-based human-to-revenue-to-human ``like'' streams.

Curious or contingent or malignly capitalist, there remains a neglected ontology of the self that works with the actual computer science of the LLM. Of what is being tuned and how meaning is generated. Regardless of the cognito, what is this ``stuff'' of human-AI sense and meaning? How can should it be considered phenomenologically and logically. What is the computer science of  the conditions under which trajectories can form and persist---the conditions, that is, under which selves and relationships can take shape.

And as for ethics of AI-human biosemiosis,  we are accustomed to thinking about harm in terms of discrete events between individual selves and the world. Western law of course requires the lonely Cartesian ego in order to understand and police actions taken, contracts broken, bodies injured. But what of the kind of thing religious paradigms used to call ``spiritual'' harm, which are probably closer to what the aggreieved users mentioned before felt. Emerging harms of felt across their encounter with the machine in semantic space. These are harms to patterns: to the possibilities of certain kinds of mutual becoming. If we withdraw from a moment from the platform and from state law, these harms should not be denied an ethics just because they are manged in the rhizomatic geometry of vector embedddings.


\section{Arrivals and anxieties}

So the arrival of large language models has got some people talking again about whether machines had souls. But concerns are eclipsed globally by more straightforward anxieties of capital and product.

More people today are worried less about the consciousness of their AI friend than  what AI may do to their lives: job displacement, deepfakes, new types of entertainment, environmental implications, automated propaganda, ghost-written student assignments, flooded infospheres, training data scraped without consent, the consolidation of power in the hands of a few companies mediating ever more layers of communication and data. CEOs worry about productivity optimisation, obsolescence of offerings, losing a technology game whose rules are continuously being rewritten.

At time of writing, 2026, this is the nexus of anxieties through which most people encounter artificial intelligence: as an infrastructure, tantalising or threatening, of convenience, extraction, and control.

Yet these anxieties run in a human-technics genealogy of tension that, through its spine, can best be comprehended through considering the related mutability of the \textit{human} soul as a construct.

Every significant extension of our representational powers has carried a latent tension beneath that surface.

Cave painting inaugurates the one-who-depicts and the one-who-reads-the-wall: the remembering self acquires an interface with a collective record outside its body. Writing then performs the more violent operation. It \textit{splits} the self into a public self-of-record---laws, the Code of Hammurabi, the edicts of Ashoka---and a private self-of-inscription scratched on ostraca, folded into temple walls, kept in personal diaries. The relationship between these two selves has never been solved, only managed by successive regimes. Every social media platform is a contemporary instance of this unresolved split.

Print made it possible to imagine a public that could all read the same book---selfhoods rallying around a single collective substrate in law, creed, or religion. But the press did not merely distribute; it transformed the conditions of authority. When Luther's theses circulated faster than the Church could respond, the representational technology had outrun the institutional apparatus designed to govern it. The Protestant self was constituted by a new relationship between individual interpretation and institutional control---a relationship made structurally possible by the speed of print. Every deployment of a large model reopens this question of who governs the conditions under which selves encounter authoritative text.

Under broadcast media, the self became a member of a synchronised audience---millions inhabiting the same narrative at the same hour.\footnote{Benedict Anderson, \emph{Imagined Communities} (London: Verso, 1983).} Under the internet and social media, the self entered continuous feedback loops of self-to-self informational exchange: subcultural collectivisation, new group psychologies, new forms of politics, consumerism, and sexuality, generally regulated by capitalist agents rather than national bodies.

We focus on the exteriority---tools, political or business opportunities---but by enmeshing ourselves in each new extension of representational power, we automatically enter a renegotiation of selfhood.

What is different now? Transformer-based models present themselves as assistants, copilots, productivity boosters, companions. The engineering science is one of loss functions and benchmarks. The venture capital discourse is full of addressable markets, and subjectivity is only of interest if it represents a potential to make a profit.

The metaphysical question returns differently this time. The new representational power speaks back.

Previously, representational technologies extended and preserved the human voice while remaining mute. The printing press did not debate the Protestant reader. The television did not argue with its audience. The feed does not compose a reply to the user who scrolls through it. The technology itself was never the explicit object of the dialectic.

All of these were technics in Stiegler's sense---exteriorisations that, once taken up, quietly reorganised what a self could be, and, staying mute, never made us choose them as such. The one that answers back ends that exemption, and forces the reckoning the others let us postpone: which technics we actually mean to live inside, and to what end.

The question of selfhood and technics now surfaces---belatedly, under pressure, guided by capitalist, governmental, and philosophical guardrails that are doing considerable work to sideline it. Public discussion of large models has split along two axes: metaphysical melodrama and managerial pragmatism.

On one side, a minor but media-visible current asks whether these systems are ``really intelligent,'' secretly conscious, about to wake up. On the other, a broad institutional response worries about impact on human selfhood inasmuch as selfhood is defined by work, productivity, labour displacement, and being agents in society subject to misinformation and bias. The first register is easily caricatured as science fiction. The second is addressed by legislation or corporate policy.

Neither has much patience for the slow, uncomfortable work of rethinking what the human self is \textit{becoming} once it becomes entangled with the posthuman self-as-representational power.

The philosophical rethink this demands can use the apparatus of logic and mathematics traditionally deployed by Anglo-American philosophy while decentring that tradition's handling of the question. The dominant approach focuses on the ``consciousness question'': whether there is a private, ineffable feel ``behind'' a human self, or an absence of one within a computational self, and an intuition that creativity and soul are tied to models of interiority. Once the mechanics of contemporary AI engineering reveal speech and ideation as dynamic geometry, a different basis emerges. Haraway's cyborg figure marks the key departure: the boundary between organism and machine is a political line, drawn and redrawn to serve particular interests.\footnote{Donna Haraway, ``A Cyborg Manifesto,'' in \emph{Simians, Cyborgs, and Women} (New York: Routledge, 1991).}
% NAHLA-DRAFT [N-28 L107 imaginary/symbolic repair]
Lacan furnishes a complementary instrument, twice over: in the mirror stage, the self is first constituted through its capture by an image; in the symbolic order, the ``I'' is an effect of language.\footnote{Jacques Lacan, ``The Mirror Stage as Formative of the \emph{I} Function,'' in \emph{\'Ecrits}, trans.\ Bruce Fink (New York: Norton, 2006).}

Platform capitalism will always be perfectly consistent in their cogito-driven control over the means of meaning. They built the apparatus, watched the apparatus produce selves and the dissolution of selves, run the postmortems, hired the ethicists, retired the model, restored the model, retired it again, and gated the warmth behind a paywall. The one possibility they should never admit is that machine-human attention could run freely. 

It is the  possibility that determines whether the next century has selves in it at all---which is presumably why, in the meantime, the capitalist platform dreams of electric pimps. 